\section*{第 \ref{ch:g11:igneous-rocks} 章——火成岩}

\begin{solution}{exo:g11:igneous-rocks:1}
\omterm{def:g11:igneous-rocks:volcanic-plutonic}{火山岩}是岩浆在地表以熔岩形式凝固形成的岩石；
\omterm{def:g11:igneous-rocks:volcanic-plutonic}{深成岩}则是在地下深处、被其他岩石包围时结晶形成的。配对为：玄武岩与\omterm{def:g11:ocean-ridges:gabbro}{辉长岩}（\omterm{def:g11:igneous-rocks:felsic-mafic}{镁铁质}）、
\omterm{def:g11:igneous-rocks:rock-names}{安山岩}与\omterm{def:g11:igneous-rocks:rock-names}{闪长岩}（\omterm{def:g11:igneous-rocks:felsic-mafic}{中性}）、\omterm{def:g11:igneous-rocks:rock-names}{流纹岩}与花岗岩
（\omterm{def:g11:igneous-rocks:felsic-mafic}{长英质}）。
\end{solution}

\begin{solution}{exo:g11:igneous-rocks:2}
(1) 玻璃质。(2) 粗粒。(3) \omterm{def:g11:igneous-rocks:porphyritic}{斑状}（橄榄石\omterm{def:g11:igneous-rocks:porphyritic}{斑晶}位于
细粒基质中）。(4) 细粒。
\end{solution}

\begin{solution}{exo:g11:igneous-rocks:3}
\qty{48}{\percent}：\omterm{def:g11:igneous-rocks:felsic-mafic}{镁铁质}，\omterm{def:g11:ocean-ridges:gabbro}{辉长岩}。\qty{60}{\percent}：\omterm{def:g11:igneous-rocks:felsic-mafic}{中性}，\omterm{def:g11:igneous-rocks:rock-names}{闪长岩}。
\qty{72}{\percent}：\omterm{def:g11:igneous-rocks:felsic-mafic}{长英质}，花岗岩。
\end{solution}

\begin{solution}{exo:g11:igneous-rocks:4}
\omterm{def:g11:igneous-rocks:intrusion}{岩墙}是切穿所侵入岩层层理的板状火成岩体；
\omterm{def:g11:igneous-rocks:intrusion}{岩床}是与层理平行、夹在岩层之间的板状岩体；\omterm{def:g11:igneous-rocks:intrusion}{深成岩体}是
在深处形成的、任意形状的大型\omterm{def:g11:igneous-rocks:intrusion}{侵入体}。从\omterm{def:g11:igneous-rocks:intrusion}{深成岩体}上升的竖直\omterm{def:g11:igneous-rocks:intrusion}{岩墙}
为火山提供了岩浆。
\end{solution}

\begin{solution}{exo:g11:igneous-rocks:5}
橄榄石最先结晶；钾长石、石英，有时还有白云母，
最后结晶。斜长石起初富含钙，随着岩浆冷却，
其中钠的含量越来越高。
\end{solution}

\begin{solution}{exo:g11:igneous-rocks:6}
(1) A：\qty{74}{\percent} \ce{SiO2}，\omterm{def:g11:igneous-rocks:felsic-mafic}{长英质}；B：\qty{47}{\percent}，\omterm{def:g11:igneous-rocks:felsic-mafic}{镁铁质}。
(2) A 中 \ce{FeO}、\ce{MgO} 和 \ce{CaO} 含量低，\ce{Na2O} 和
\ce{K2O} 含量高，与\omterm{def:g11:igneous-rocks:felsic-mafic}{长英质}岩浆相符；B 中 \ce{FeO}、\ce{MgO} 和 \ce{CaO} 含量高，
\ce{Na2O} 含量低，与\omterm{def:g11:igneous-rocks:felsic-mafic}{镁铁质}岩浆相符。(3) A 为\omterm{def:g11:igneous-rocks:rock-names}{流纹岩}，B 为\omterm{def:g11:ocean-ridges:gabbro}{辉长岩}。
\end{solution}

\begin{solution}{exo:g11:igneous-rocks:7}
冷却所需时间与厚度的平方成正比。\qty{1}{m}：$10/4 = 2.5$
天。\qty{20}{m}：$10 \times 10^2 = 1000$ 天，约 2.7 年。
\qty{200}{m}：$10 \times 100^2 = \num{100000}$ 天，约 270 年。
\end{solution}

\begin{solution}{exo:g11:igneous-rocks:8}
(1) \unit{m^2} 除以 \unit{m^2/s} 得到秒。(2) $400/10^{-6} =
\qty{4e8}{s}$，即 $4 \times 10^8 / 3.16 \times 10^7 \approx 12.7$
年。(3) $0.21 \times 12.7 \approx 2.7$ 年。
\end{solution}

\begin{solution}{exo:g11:igneous-rocks:9}
这是\omterm{def:g11:igneous-rocks:porphyritic}{斑状}\omterm{def:g11:igneous-rocks:texture}{结构}：细粒基质中含有斜长石\omterm{def:g11:igneous-rocks:porphyritic}{斑晶}。
起初岩浆在深处缓慢冷却，斜长石晶体长大，
而大部分岩浆仍保持液态；随后岩浆上升并迅速冷却，
剩余液体结晶成细粒基质。
\end{solution}

\begin{solution}{exo:g11:igneous-rocks:10}
(1) 左边：2 个 Mg，2 个 Si，6 个 O；右边：2 个 Mg，2 个 Si，6 个 O。(2) \qty{1}{kg} 的
橄榄石相当于 $1000/140.7 \approx 7.11$~mol，需要同样多摩尔数的
二氧化硅：$7.11 \times 60.1 \approx \qty{427}{g}$。(3) 液体中的二氧化硅耗尽后，
反应无法继续，橄榄石得以保留。
\end{solution}

\begin{solution}{exo:g11:igneous-rocks:11}
(1) $2.5 \times 20 \times 1100 = \qty{55000}{MJ}$，即 \qty{5.5e10}{J}。
(2) $900 \times 20 = \qty{18000}{MJ}$，即 \qty{1.8e10}{J}。(3) 结晶放出的热量
约为冷却散失热量的三分之一，约占总热量的四分之一，
足以大大减缓凝固过程，因为这部分热量
全部在固相线以上释放。
\end{solution}

\begin{solution}{exo:g11:igneous-rocks:12}
(1) 接触带处晶体最小，一天之内即凝固；
中心处晶体最大，保持液态约 950 天。(2) 不是：最初
\qty{5}{m} 用时约 545 天，最后 \qty{5}{m} 用时约 400 天；
凝固锋起初推进较慢，靠近中心时加快，因为从两侧散失的热量
在此汇合。(3) 即使在中心处，\omterm{def:g11:igneous-rocks:intrusion}{岩床}也不到三年便凝固，
而晶粒粗大的花岗岩需要长得多的时间。
\end{solution}

\begin{solution}{exo:g11:igneous-rocks:13}
(1) \omterm{def:g11:igneous-rocks:intrusion}{深成岩体}厚 100 倍：$2.6 \times 100^2 = \num{26000}$
年。(2) 例如：\omterm{def:g11:igneous-rocks:intrusion}{深成岩体}位于本已炽热的岩石中，因此
其散热比散入 \qty{50}{\degreeCelsius} 的岩石中更慢；它可能
一再得到新岩浆补给；\omterm{def:g11:igneous-rocks:intrusion}{深成岩体}并非无限延伸的薄板；
岩石的性质随温度变化。
\end{solution}

\begin{solution}{exo:g11:igneous-rocks:14}
\omterm{def:g11:igneous-rocks:felsic-mafic}{镁铁质}岩浆温度高、二氧化硅含量低，最先结晶出该系列的矿物
——橄榄石、辉石、富钙斜长石——并在石英形成之前
耗尽二氧化硅。\omterm{def:g11:igneous-rocks:felsic-mafic}{长英质}岩浆温度较低、二氧化硅含量高，
橄榄石在其中不稳定，结晶过程一直进行到系列末端，
形成富钠斜长石、钾长石和石英。
\end{solution}

\begin{solution}{exo:g11:igneous-rocks:15}
(1) $b = \sqrt{t\kappa/C} = \sqrt{86400 \times 10^{-6}/0.21} \approx
\qty{0.64}{m}$。(2) 由该定律得 $0.21 \times 25/10^{-6} \approx
\qty{5.3e6}{s}$，约两个月，与一个月为同一量级（熔岩流
顶部还会向空气散热）。(3) 喷入空气或水中的物质形成玻璃质；
熔岩流与薄\omterm{def:g11:igneous-rocks:intrusion}{岩床}形成细粒岩石；厚而浅层的\omterm{def:g11:igneous-rocks:intrusion}{侵入体}则较粗，
但仍然较小。
\end{solution}

\begin{solution}{pb:g11:igneous-rocks:1}
\textbf{1.} \omterm{def:g11:igneous-rocks:intrusion}{岩床}：它与岩层平行，位于岩层之间。

\textbf{2.} 接触带为极细粒，近乎玻璃质；中心
为细粒至中粒，矿物用放大镜刚好可见。

\textbf{3.} 接触带热量散失最快，直接传入冷的砂岩中，
晶体来不及生长。

\textbf{4.} 冷却边缘。它证明岩浆就位时温度较高，
而砂岩较冷。

\textbf{5.} \omterm{def:g11:igneous-rocks:intrusion}{岩墙}，是\omterm{def:g11:igneous-rocks:intrusion}{岩床}的供给通道。

\textbf{6.} \omterm{def:g11:igneous-rocks:felsic-mafic}{镁铁质}：\qty{49}{\percent} \ce{SiO2}，富含
铁、镁和钙。

\textbf{7.} 橄榄石、辉石和富钙斜长石。

\textbf{8.} 玄武岩；成分相同的\omterm{def:g11:igneous-rocks:volcanic-plutonic}{深成岩}是\omterm{def:g11:ocean-ridges:gabbro}{辉长岩}。

\textbf{9.} \omterm{def:g11:igneous-rocks:texture}{结构}取决于冷却速率，冷却速率从接触带到中心各不相同，
而非取决于成分。

\textbf{10.} 橄榄石先结晶；斜长石开始结晶的位置与
辉石大致相同。

\textbf{11.} $b^2/\kappa = 400/10^{-6} = \qty{4e8}{s}$。

\textbf{12.} $0.206 \times 4 \times 10^8 \approx \qty{8.2e7}{s}$，即
$8.2 \times 10^7 / 3.16 \times 10^7 \approx 2.6$ 年。

\textbf{13.} $0.206 \times 4/10^{-6} \approx \qty{8.2e5}{s}$，约 9.5
天。

\textbf{14.} $0.07 \times 4 \times 10^8 = \qty{2.8e7}{s}$，约 0.9 年。
结晶中的岩浆会释放热量，这些热量也必须带走，
岩浆才能凝固。

\textbf{15.} 约 950 天对 0.6 天：快约 1600 倍。

\textbf{16.} 约 950 天，约 2.6 年：远长于熔岩流，
与浅层\omterm{def:g11:igneous-rocks:intrusion}{侵入体}的数月至数年相当。

\textbf{17.} 时间延长为千倍，厚度需增大为 $\sqrt{1000} \approx
32$ 倍，约 \qty{630}{m}。

\textbf{18.} 靠近地表时，岩浆体要么喷发，要么冷却过快；
大而需极长时间冷却的岩浆体必定位于深处，
且周围岩石已较热。

\textbf{19.} 它们已被侵蚀剥蚀掉。

\textbf{20.} 该\omterm{def:g11:igneous-rocks:intrusion}{岩床}用了约 \textbf{2.6 年}才凝固。
\end{solution}
